\begin{abstract}
Federated learning lets many clients train one model without pooling their data, yet the updates they share can
still expose individual records, so deployments add differential privacy (DP). DP does not cost every client the
same: clients whose data the model fits poorly lose the most accuracy, which turns a privacy guarantee into a
fairness problem. Personalised federated learning is known to soften this cost, but its main mechanisms---local
fine-tuning, meta-learning, client clustering and parameter decoupling---have been tested in separate studies,
against different baselines and mostly without DP. We compare one representative of each family under a single
record-level privacy budget. Every client trains with per-example DP-SGD and answers the same private queries, and
the noise of each run is calibrated so that all mechanisms spend exactly the same $(\eps,\delta)$, including the
cost of the loss reports needed by clustering and by the q-FFL fairness objective. On CIFAR-10 split over 20
clients with a Dirichlet label partition, we sweep $\eps$ from 0.3 to 8 and measure mean and worst-10\% client
accuracy, Jain's index, the gain over the DP global model, membership-inference leakage and communication cost;
on the Adult census data we add demographic-parity and equal-opportunity gaps.
At $\eps=2$, FedPer roughly doubles DP-FedAvg's mean client accuracy (54.1\% versus 27.2\%) and triples its
worst-10\% accuracy, whereas the q-FFL fairness objective stalls under DP.
The outcome is a recommended operating configuration: which mechanism to use at a given privacy budget.
\end{abstract}

\begin{IEEEkeywords}
Differential privacy, fairness, federated learning, membership inference, non-IID data, personalisation.
\end{IEEEkeywords}
